\documentclass[10pt]{article}
% Hojas resumen: un A4 por delante y por detrás, dos columnas por cara.
% Uso: \documentclass[10pt]{article} % Hojas resumen: un A4 por delante y por detrás, dos columnas por cara.
% Uso: \documentclass[10pt]{article} % Hojas resumen: un A4 por delante y por detrás, dos columnas por cara.
% Uso: \documentclass[10pt]{article} \input{estilo-resumen.tex} \bloque{analisis|algebra|azar} \cabecera{n}{Título}{Lecciones}
\usepackage[utf8]{inputenc}
\usepackage[T1]{fontenc}
\usepackage[spanish,es-nodecimaldot]{babel}
\usepackage{lmodern}
\usepackage{amsmath,amssymb}
\usepackage{graphicx}
\usepackage[a4paper,left=10mm,right=10mm,top=15mm,bottom=13mm,headheight=11pt,headsep=3mm,footskip=6mm]{geometry}
\usepackage{multicol}
\setlength{\columnsep}{6mm}
\setlength{\columnseprule}{0.3pt}
\usepackage{xcolor}
\usepackage{enumitem}
\setlist{nosep,leftmargin=1.1em,itemsep=1pt}
\usepackage{array}
\usepackage{needspace}
\usepackage{pgfplots}
\usepgfplotslibrary{fillbetween}
\pgfplotsset{compat=1.18,
  mini/.style={width=4.3cm,height=3.3cm,scale only axis=false,axis lines=middle,xtick=\empty,ytick=\empty,samples=120,
    xlabel={\scriptsize$x$},ylabel={\scriptsize$y$},title style={font=\scriptsize,yshift=-4pt},clip=true,
    every axis plot/.append style={thick,color=acento}}}
\usepackage{ragged2e}
\usepackage[most]{tcolorbox}
\usepackage{fancyhdr,lastpage}
\usepackage[hidelinks,pdfencoding=auto]{hyperref}
\hypersetup{pdflang=es-ES,pdfsubject={Matemáticas Aplicadas a las Ciencias Sociales II}}

\definecolor{analisis}{HTML}{1D4ED8}
\definecolor{algebra}{HTML}{86198F}
\definecolor{azar}{HTML}{9A3412}
\definecolor{aviso}{HTML}{B91C1C}
\colorlet{acento}{analisis}
\newcommand{\bloque}[1]{\colorlet{acento}{#1}}

\setlength{\parindent}{0pt}
\setlength{\parskip}{2pt}
\renewcommand{\arraystretch}{1.2}
\AtBeginDocument{\RaggedRight\fontsize{9.8}{11.9}\selectfont\setlength{\abovedisplayskip}{3pt}\setlength{\belowdisplayskip}{3pt}\setlength{\abovedisplayshortskip}{1pt}\setlength{\belowdisplayshortskip}{1pt}}

\newcommand{\temanum}{}\newcommand{\tematitulo}{}
\pagestyle{fancy}
\fancyhf{}
\fancyhead[L]{\footnotesize\color{acento}\textbf{Hoja resumen}\quad Tema \temanum. \tematitulo}
\fancyhead[R]{\footnotesize Matemáticas Aplicadas a las Ciencias Sociales II --- 2.\textsuperscript{o} Bachillerato}
\fancyfoot[L]{\scriptsize IES Al-Ándalus --- Juan María Gámez Ortiz, con la ayuda de Claude y de Codex --- CC BY-NC 4.0}
\fancyfoot[R]{\scriptsize Cara \thepage\ de \pageref{LastPage}}
\renewcommand{\headrulewidth}{0.4pt}

% Cabecera de la primera cara: número, título y lecciones del tema.
\newcommand{\cabecera}[3]{%
  \renewcommand{\temanum}{#1}\renewcommand{\tematitulo}{#2}%
  \hypersetup{pdftitle={Hoja resumen. Tema #1. #2}}%
  \noindent\begin{tcolorbox}[enhanced,colback=acento!8,colframe=acento,boxrule=0.6pt,arc=2pt,left=4pt,right=4pt,top=2pt,bottom=2pt]
  \begin{minipage}[c]{1.1cm}\includegraphics[width=1cm]{logo-al-andalus.png}\end{minipage}%
  \begin{minipage}[c]{\dimexpr\linewidth-1.2cm}{\large\bfseries\color{acento}Tema #1. #2}\\[1pt]{\scriptsize #3}\end{minipage}
  \end{tcolorbox}\vspace{-2pt}}

% Apartado con filete del color del bloque.
\newcommand{\seccion}[1]{\par\needspace{5\baselineskip}\vspace{3pt}{\color{acento}\bfseries\normalsize #1}\par\vspace{-1pt}{\color{acento}\hrule height 0.5pt}\vspace{2pt}}
\newcommand{\sub}[1]{\par\needspace{3\baselineskip}\vspace{1pt}\textbf{\color{acento}#1}\enspace}

% Fórmula clave, aviso de errores y ejemplo.
\tcbset{before upper=\RaggedRight}
\newtcolorbox{formula}{enhanced,colback=acento!6,colframe=acento!70,boxrule=0.5pt,arc=1.5pt,left=3pt,right=3pt,top=0pt,bottom=0pt,before skip=2pt,after skip=2pt}
\newtcolorbox{ojo}{enhanced,colback=aviso!5,colframe=aviso!60,boxrule=0.5pt,arc=1.5pt,left=3pt,right=3pt,top=1pt,bottom=1pt,before skip=3pt,after skip=2pt,
  fonttitle=\bfseries\footnotesize,coltitle=aviso,colbacktitle=aviso!5,title={Ojo},attach title to upper={\enspace}}
\newtcolorbox{ejemplo}{enhanced,colback=white,colframe=acento!45,boxrule=0.4pt,arc=1.5pt,left=3pt,right=3pt,top=1pt,bottom=1pt,before skip=3pt,after skip=2pt,
  fonttitle=\bfseries\footnotesize,coltitle=acento,colbacktitle=white,title={Ejemplo},attach title to upper={\enspace}}
\newcommand{\dx}{\,dx}
 \bloque{analisis|algebra|azar} \cabecera{n}{Título}{Lecciones}
\usepackage[utf8]{inputenc}
\usepackage[T1]{fontenc}
\usepackage[spanish,es-nodecimaldot]{babel}
\usepackage{lmodern}
\usepackage{amsmath,amssymb}
\usepackage{graphicx}
\usepackage[a4paper,left=10mm,right=10mm,top=15mm,bottom=13mm,headheight=11pt,headsep=3mm,footskip=6mm]{geometry}
\usepackage{multicol}
\setlength{\columnsep}{6mm}
\setlength{\columnseprule}{0.3pt}
\usepackage{xcolor}
\usepackage{enumitem}
\setlist{nosep,leftmargin=1.1em,itemsep=1pt}
\usepackage{array}
\usepackage{needspace}
\usepackage{pgfplots}
\usepgfplotslibrary{fillbetween}
\pgfplotsset{compat=1.18,
  mini/.style={width=4.3cm,height=3.3cm,scale only axis=false,axis lines=middle,xtick=\empty,ytick=\empty,samples=120,
    xlabel={\scriptsize$x$},ylabel={\scriptsize$y$},title style={font=\scriptsize,yshift=-4pt},clip=true,
    every axis plot/.append style={thick,color=acento}}}
\usepackage{ragged2e}
\usepackage[most]{tcolorbox}
\usepackage{fancyhdr,lastpage}
\usepackage[hidelinks,pdfencoding=auto]{hyperref}
\hypersetup{pdflang=es-ES,pdfsubject={Matemáticas Aplicadas a las Ciencias Sociales II}}

\definecolor{analisis}{HTML}{1D4ED8}
\definecolor{algebra}{HTML}{86198F}
\definecolor{azar}{HTML}{9A3412}
\definecolor{aviso}{HTML}{B91C1C}
\colorlet{acento}{analisis}
\newcommand{\bloque}[1]{\colorlet{acento}{#1}}

\setlength{\parindent}{0pt}
\setlength{\parskip}{2pt}
\renewcommand{\arraystretch}{1.2}
\AtBeginDocument{\RaggedRight\fontsize{9.8}{11.9}\selectfont\setlength{\abovedisplayskip}{3pt}\setlength{\belowdisplayskip}{3pt}\setlength{\abovedisplayshortskip}{1pt}\setlength{\belowdisplayshortskip}{1pt}}

\newcommand{\temanum}{}\newcommand{\tematitulo}{}
\pagestyle{fancy}
\fancyhf{}
\fancyhead[L]{\footnotesize\color{acento}\textbf{Hoja resumen}\quad Tema \temanum. \tematitulo}
\fancyhead[R]{\footnotesize Matemáticas Aplicadas a las Ciencias Sociales II --- 2.\textsuperscript{o} Bachillerato}
\fancyfoot[L]{\scriptsize IES Al-Ándalus --- Juan María Gámez Ortiz, con la ayuda de Claude y de Codex --- CC BY-NC 4.0}
\fancyfoot[R]{\scriptsize Cara \thepage\ de \pageref{LastPage}}
\renewcommand{\headrulewidth}{0.4pt}

% Cabecera de la primera cara: número, título y lecciones del tema.
\newcommand{\cabecera}[3]{%
  \renewcommand{\temanum}{#1}\renewcommand{\tematitulo}{#2}%
  \hypersetup{pdftitle={Hoja resumen. Tema #1. #2}}%
  \noindent\begin{tcolorbox}[enhanced,colback=acento!8,colframe=acento,boxrule=0.6pt,arc=2pt,left=4pt,right=4pt,top=2pt,bottom=2pt]
  \begin{minipage}[c]{1.1cm}\includegraphics[width=1cm]{logo-al-andalus.png}\end{minipage}%
  \begin{minipage}[c]{\dimexpr\linewidth-1.2cm}{\large\bfseries\color{acento}Tema #1. #2}\\[1pt]{\scriptsize #3}\end{minipage}
  \end{tcolorbox}\vspace{-2pt}}

% Apartado con filete del color del bloque.
\newcommand{\seccion}[1]{\par\needspace{5\baselineskip}\vspace{3pt}{\color{acento}\bfseries\normalsize #1}\par\vspace{-1pt}{\color{acento}\hrule height 0.5pt}\vspace{2pt}}
\newcommand{\sub}[1]{\par\needspace{3\baselineskip}\vspace{1pt}\textbf{\color{acento}#1}\enspace}

% Fórmula clave, aviso de errores y ejemplo.
\tcbset{before upper=\RaggedRight}
\newtcolorbox{formula}{enhanced,colback=acento!6,colframe=acento!70,boxrule=0.5pt,arc=1.5pt,left=3pt,right=3pt,top=0pt,bottom=0pt,before skip=2pt,after skip=2pt}
\newtcolorbox{ojo}{enhanced,colback=aviso!5,colframe=aviso!60,boxrule=0.5pt,arc=1.5pt,left=3pt,right=3pt,top=1pt,bottom=1pt,before skip=3pt,after skip=2pt,
  fonttitle=\bfseries\footnotesize,coltitle=aviso,colbacktitle=aviso!5,title={Ojo},attach title to upper={\enspace}}
\newtcolorbox{ejemplo}{enhanced,colback=white,colframe=acento!45,boxrule=0.4pt,arc=1.5pt,left=3pt,right=3pt,top=1pt,bottom=1pt,before skip=3pt,after skip=2pt,
  fonttitle=\bfseries\footnotesize,coltitle=acento,colbacktitle=white,title={Ejemplo},attach title to upper={\enspace}}
\newcommand{\dx}{\,dx}
 \bloque{analisis|algebra|azar} \cabecera{n}{Título}{Lecciones}
\usepackage[utf8]{inputenc}
\usepackage[T1]{fontenc}
\usepackage[spanish,es-nodecimaldot]{babel}
\usepackage{lmodern}
\usepackage{amsmath,amssymb}
\usepackage{graphicx}
\usepackage[a4paper,left=10mm,right=10mm,top=15mm,bottom=13mm,headheight=11pt,headsep=3mm,footskip=6mm]{geometry}
\usepackage{multicol}
\setlength{\columnsep}{6mm}
\setlength{\columnseprule}{0.3pt}
\usepackage{xcolor}
\usepackage{enumitem}
\setlist{nosep,leftmargin=1.1em,itemsep=1pt}
\usepackage{array}
\usepackage{needspace}
\usepackage{pgfplots}
\usepgfplotslibrary{fillbetween}
\pgfplotsset{compat=1.18,
  mini/.style={width=4.3cm,height=3.3cm,scale only axis=false,axis lines=middle,xtick=\empty,ytick=\empty,samples=120,
    xlabel={\scriptsize$x$},ylabel={\scriptsize$y$},title style={font=\scriptsize,yshift=-4pt},clip=true,
    every axis plot/.append style={thick,color=acento}}}
\usepackage{ragged2e}
\usepackage[most]{tcolorbox}
\usepackage{fancyhdr,lastpage}
\usepackage[hidelinks,pdfencoding=auto]{hyperref}
\hypersetup{pdflang=es-ES,pdfsubject={Matemáticas Aplicadas a las Ciencias Sociales II}}

\definecolor{analisis}{HTML}{1D4ED8}
\definecolor{algebra}{HTML}{86198F}
\definecolor{azar}{HTML}{9A3412}
\definecolor{aviso}{HTML}{B91C1C}
\colorlet{acento}{analisis}
\newcommand{\bloque}[1]{\colorlet{acento}{#1}}

\setlength{\parindent}{0pt}
\setlength{\parskip}{2pt}
\renewcommand{\arraystretch}{1.2}
\AtBeginDocument{\RaggedRight\fontsize{9.8}{11.9}\selectfont\setlength{\abovedisplayskip}{3pt}\setlength{\belowdisplayskip}{3pt}\setlength{\abovedisplayshortskip}{1pt}\setlength{\belowdisplayshortskip}{1pt}}

\newcommand{\temanum}{}\newcommand{\tematitulo}{}
\pagestyle{fancy}
\fancyhf{}
\fancyhead[L]{\footnotesize\color{acento}\textbf{Hoja resumen}\quad Tema \temanum. \tematitulo}
\fancyhead[R]{\footnotesize Matemáticas Aplicadas a las Ciencias Sociales II --- 2.\textsuperscript{o} Bachillerato}
\fancyfoot[L]{\scriptsize IES Al-Ándalus --- Juan María Gámez Ortiz, con la ayuda de Claude y de Codex --- CC BY-NC 4.0}
\fancyfoot[R]{\scriptsize Cara \thepage\ de \pageref{LastPage}}
\renewcommand{\headrulewidth}{0.4pt}

% Cabecera de la primera cara: número, título y lecciones del tema.
\newcommand{\cabecera}[3]{%
  \renewcommand{\temanum}{#1}\renewcommand{\tematitulo}{#2}%
  \hypersetup{pdftitle={Hoja resumen. Tema #1. #2}}%
  \noindent\begin{tcolorbox}[enhanced,colback=acento!8,colframe=acento,boxrule=0.6pt,arc=2pt,left=4pt,right=4pt,top=2pt,bottom=2pt]
  \begin{minipage}[c]{1.1cm}\includegraphics[width=1cm]{logo-al-andalus.png}\end{minipage}%
  \begin{minipage}[c]{\dimexpr\linewidth-1.2cm}{\large\bfseries\color{acento}Tema #1. #2}\\[1pt]{\scriptsize #3}\end{minipage}
  \end{tcolorbox}\vspace{-2pt}}

% Apartado con filete del color del bloque.
\newcommand{\seccion}[1]{\par\needspace{5\baselineskip}\vspace{3pt}{\color{acento}\bfseries\normalsize #1}\par\vspace{-1pt}{\color{acento}\hrule height 0.5pt}\vspace{2pt}}
\newcommand{\sub}[1]{\par\needspace{3\baselineskip}\vspace{1pt}\textbf{\color{acento}#1}\enspace}

% Fórmula clave, aviso de errores y ejemplo.
\tcbset{before upper=\RaggedRight}
\newtcolorbox{formula}{enhanced,colback=acento!6,colframe=acento!70,boxrule=0.5pt,arc=1.5pt,left=3pt,right=3pt,top=0pt,bottom=0pt,before skip=2pt,after skip=2pt}
\newtcolorbox{ojo}{enhanced,colback=aviso!5,colframe=aviso!60,boxrule=0.5pt,arc=1.5pt,left=3pt,right=3pt,top=1pt,bottom=1pt,before skip=3pt,after skip=2pt,
  fonttitle=\bfseries\footnotesize,coltitle=aviso,colbacktitle=aviso!5,title={Ojo},attach title to upper={\enspace}}
\newtcolorbox{ejemplo}{enhanced,colback=white,colframe=acento!45,boxrule=0.4pt,arc=1.5pt,left=3pt,right=3pt,top=1pt,bottom=1pt,before skip=3pt,after skip=2pt,
  fonttitle=\bfseries\footnotesize,coltitle=acento,colbacktitle=white,title={Ejemplo},attach title to upper={\enspace}}
\newcommand{\dx}{\,dx}

\bloque{analisis}
\begin{document}
\cabecera{2}{Representación de funciones}{Lecciones 2.1--2.9: esquema general, polinómicas, racionales, exponenciales, logarítmicas, funciones a trozos y valor absoluto, traslaciones, expansión y contracción, optimización.}
\begin{multicols}{2}

\seccion{1. Esquema general}
\begin{enumerate}
  \item \textbf{Dominio:} dónde existe la fórmula (denominadores \(\ne0\), radicandos de índice par \(\ge0\), argumentos del logaritmo \(>0\)).
  \item \textbf{Cortes, simetría y signo:} con el eje \(Y\), \(x=0\); con el eje \(X\), \(f(x)=0\). Par si \(f(-x)=f(x)\) (simétrica respecto al eje \(Y\)); impar si \(f(-x)=-f(x)\) (respecto al origen).
  \item \textbf{Continuidad y límites} en los bordes del dominio.
  \item \textbf{Asíntotas:} vertical, horizontal u oblicua.
  \item \textbf{Derivada y monotonía:} signo de \(f'\).
  \item \textbf{Extremos relativos.}
  \item \textbf{Tabla y boceto} con toda la información.
\end{enumerate}

\seccion{2. Monotonía y extremos}
\begin{formula}
\(f'(x)>0\) en un intervalo \(\Rightarrow f\) crece; \quad \(f'(x)<0\Rightarrow f\) decrece.
\end{formula}
\begin{itemize}
  \item \textbf{Puntos críticos:} \(f'(a)=0\) o \(f'\) no existe.
  \item \textbf{Máximo relativo:} \(f'\) pasa de \(+\) a \(-\). \textbf{Mínimo:} de \(-\) a \(+\). Si no cambia de signo, no hay extremo.
  \item Con la segunda derivada: \(f'(a)=0\) y \(f''(a)<0\Rightarrow\) máximo; \(f''(a)>0\Rightarrow\) mínimo.
  \item \textbf{Curvatura} (ampliación): \(f''>0\), forma de cuenco; \(f''<0\), forma de arco. \textbf{Inflexión}: cambia la curvatura.
\end{itemize}
\begin{ejemplo}
\(f(x)=x^3-3x\): dominio \(\mathbb R\), impar, cortes en \(x=0,\pm\sqrt3\).
\(f'(x)=3x^2-3=0\Rightarrow x=\pm1\).
\[\begin{array}{c|ccccc} x & (-\infty,-1) & -1 & (-1,1) & 1 & (1,+\infty)\\\hline f' & + & 0 & - & 0 & +\\ f & \nearrow & 2 & \searrow & -2 & \nearrow\end{array}\]
Máximo \((-1,2)\), mínimo \((1,-2)\). \(f''(x)=6x\): inflexión en \((0,0)\).
\end{ejemplo}

\seccion{3. Asíntotas}
\begin{formula}
\textbf{Vertical} \(x=a\): \(\lim_{x\to a^\pm}f(x)=\pm\infty\).\\
\textbf{Horizontal} \(y=L\): \(\lim_{x\to\pm\infty}f(x)=L\).\\
\textbf{Oblicua} \(y=mx+n\):
\[m=\lim_{x\to\pm\infty}\frac{f(x)}{x}\ (\ne0),\qquad n=\lim_{x\to\pm\infty}\bigl(f(x)-mx\bigr)\]
\end{formula}
Se estudian por separado \(+\infty\) y \(-\infty\). La gráfica puede cortar una asíntota horizontal u oblicua; nunca una vertical.

\seccion{Gráficas de referencia}
\begin{center}
\begin{tikzpicture}
\begin{axis}[mini,title={\(y=x^2-1\)},xmin=-2.2,xmax=2.2,ymin=-1.6,ymax=3.2]\addplot[domain=-2.1:2.1]{x^2-1};\end{axis}
\end{tikzpicture}\hfill
\begin{tikzpicture}
\begin{axis}[mini,title={\(y=2^x\) y \(y=0{,}5^x\)},xmin=-3,xmax=3,ymin=-0.5,ymax=6]\addplot[domain=-3:2.6]{2^x};\addplot[domain=-2.6:3,dashed]{0.5^x};\end{axis}
\end{tikzpicture}\\[2pt]
\begin{tikzpicture}
\begin{axis}[mini,title={\(y=\ln x\)},xmin=-0.6,xmax=6,ymin=-3,ymax=2.2]\addplot[domain=0.05:6]{ln(x)};\end{axis}
\end{tikzpicture}\hfill
\begin{tikzpicture}
\begin{axis}[mini,title={\(y=\frac1x\)},xmin=-3,xmax=3,ymin=-3,ymax=3]\addplot[domain=-3:-0.34]{1/x};\addplot[domain=0.34:3]{1/x};\end{axis}
\end{tikzpicture}
\end{center}
\begin{tabular}{@{}p{1.6cm}p{1.6cm}p{4.6cm}@{}}
\textbf{Familia} & \textbf{Dominio} & \textbf{Asíntotas y forma}\\ \hline
Polinómica & \(\mathbb R\) & Sin asíntotas; ramas por \(a_nx^n\)\\
Racional & \(Q(x)\ne0\) & Verticales, horizontal u oblicua\\
\(a^x\) & \(\mathbb R\) & \(y=0\); siempre positiva\\
\(\log_ax\) & \(x>0\) & \(x=0\); pasa por \((1,0)\)\\
\(|g(x)|\) & el de \(g\) & Pico donde \(g\) cambia de signo
\end{tabular}

\seccion{4. Polinómicas}
Dominio \(\mathbb R\), continuas y sin asíntotas. Un polinomio de grado \(n\) tiene como mucho \(n\) raíces y \(n-1\) extremos. Las ramas en \(\pm\infty\) las decide el término de mayor grado: \(a_nx^n\).
\sub{Raíces y signo.} Una raíz simple cruza el eje; una doble lo toca sin cruzarlo (cambia o no cambia el signo).

\seccion{5. Racionales \(f=\frac PQ\)}
\begin{formula}
\[D=\{x\in\mathbb R:\ Q(x)\ne0\}\]
\end{formula}
\begin{itemize}
  \item Si \(Q(a)=0\) y \(P(a)\ne0\): asíntota vertical \(x=a\).
  \item Si \(P(a)=Q(a)=0\): se simplifica; puede quedar un \textbf{hueco} (discontinuidad evitable) en lugar de una asíntota.
  \item Grado \(P<\) grado \(Q\): \(y=0\). Grados iguales: \(y=\frac{a_n}{b_n}\). Grado \(P=\) grado \(Q+1\): asíntota oblicua.
\end{itemize}
\begin{ejemplo}
\(f(x)=\dfrac{x^2+1}{x}=x+\dfrac1x\): AV \(x=0\), AO \(y=x\). \(f'(x)=\dfrac{x^2-1}{x^2}=0\Rightarrow x=\pm1\): máximo \((-1,-2)\) y mínimo \((1,2)\).
\end{ejemplo}

\seccion{6. Exponenciales \(y=a^x\) \((a>0,\ a\ne1)\)}
Dominio \(\mathbb R\), recorrido \((0,+\infty)\), pasa por \((0,1)\), asíntota horizontal \(y=0\). Si \(a>1\) crece; si \(0<a<1\) decrece.
\begin{formula}
\[V(t)=V_0(1+r)^t,\qquad r>-1\]
\(r>0\): crecimiento (interés compuesto). \(r<0\): depreciación.
\end{formula}
\begin{ejemplo}
Una máquina de 20\,000\,euros pierde el 15\,\% de su valor cada año: \(V(t)=20\,000\cdot0{,}85^{\,t}\). A los 3 años: \(V(3)\approx12\,282{,}50\)\,euros.
\end{ejemplo}

\seccion{7. Logarítmicas \(y=\log_ax\)}
\begin{formula}
\[\log_a x=y\iff a^y=x\]
\[\log_a(uv)=\log_a u+\log_a v,\qquad \log_a\frac uv=\log_a u-\log_a v\]
\[\log_a(u^r)=r\log_a u,\qquad \log_a x=\frac{\ln x}{\ln a}\]
\end{formula}
Dominio \((0,+\infty)\), recorrido \(\mathbb R\), pasa por \((1,0)\), asíntota vertical \(x=0\). Es la inversa de \(a^x\): gráficas simétricas respecto a \(y=x\).
\begin{ejemplo}
\(\log_2(x-1)=3\Rightarrow x-1=2^3\Rightarrow x=9\). Comprobación: \(x-1=8>0\) y \(\log_28=3\). En las ecuaciones con logaritmos hay que comprobar que cada candidato está en el dominio.
\end{ejemplo}
\begin{ejemplo}
\textquestiondown{}Cuándo se duplica un capital al 5\,\% anual? \(2=1{,}05^{\,t}\Rightarrow t=\dfrac{\ln2}{\ln1{,}05}\approx14{,}2\) años.
\end{ejemplo}

\seccion{8. A trozos y valor absoluto}
\begin{formula}
\[|x|=\begin{cases}-x,&x<0\\ x,&x\ge0\end{cases}\]
\end{formula}
Para \(|g(x)|\): se estudia el signo de \(g\) y se escribe por trozos. En cada punto de unión se comprueban los laterales: si coinciden con el valor, no hay salto; si además coinciden las derivadas laterales, no hay pico.
\begin{ejemplo}
\(g(x)=|x-1|-2\) vale \(-x-1\) si \(x<1\) y \(x-3\) si \(x\ge1\): una \guillemotleft{}V\guillemotright{} con vértice en \((1,-2)\) que corta el eje \(X\) en \(x=-1\) y \(x=3\).
\end{ejemplo}

\seccion{9. Transformaciones de la gráfica de \(f\)}
\begin{formula}
\textbf{Traslación:}
\[g(x)=f(x-a)+k:\qquad (u,\,f(u))\mapsto(u+a,\ f(u)+k)\]
\textbf{Escala:}
\[g(x)=a\,f(bx):\qquad (u,\,v)\mapsto(u/b,\ av)\]
\end{formula}
\begin{itemize}
  \item \(a>0\) en \(f(x-a)\): a la \textbf{derecha}; \(k>0\): hacia \textbf{arriba}.
  \item \(|a|>1\) estira en vertical; \(|b|>1\) comprime en horizontal.
  \item \(-f(x)\): simetría respecto al eje \(X\). \(f(-x)\): respecto al eje \(Y\).
  \item Las asíntotas y los extremos se transforman con la misma regla.
\end{itemize}
\begin{ejemplo}
\(g(x)=(x-2)^2-3\) es \(y=x^2\) trasladada 2 a la derecha y 3 abajo: vértice \((2,-3)\).
\(r(x)=\frac1{x+2}+1\): asíntotas \(x=-2\) e \(y=1\).
\end{ejemplo}

\seccion{10. Optimización}
\begin{enumerate}
  \item Elige la variable y escribe la función que se optimiza.
  \item Fija el \textbf{dominio} con las restricciones del enunciado.
  \item Deriva, iguala a 0 y obtén los candidatos interiores.
  \item Compara con los \textbf{extremos del intervalo}.
  \item Responde con unidades.
\end{enumerate}
\begin{ejemplo}
\(B(q)=-q^2+40q-300\) miles de euros, \(q\in[0,30]\) toneladas. \(B'(q)=-2q+40=0\Rightarrow q=20\), \(B(20)=100\). Si la capacidad fuera \(q\le15\), \(B\) crece en \([0,15]\) y el máximo es \(B(15)=75\).
\end{ejemplo}
\begin{ejemplo}
\textbf{Coste medio.} \(C(q)=q^2+100\) euros por \(q\) unidades. \(\overline C(q)=\dfrac{C(q)}q=q+\dfrac{100}q\), \(q>0\).
\(\overline C'(q)=1-\dfrac{100}{q^2}=0\Rightarrow q=10\); \(\overline C'\) pasa de \(-\) a \(+\): mínimo, \(\overline C(10)=20\)\,euros/unidad.
\end{ejemplo}
\begin{ejemplo}
\textbf{Geometría.} Rectángulo de perímetro 40 m y área máxima: lados \(x\) y \(20-x\), \(A(x)=x(20-x)\), \(x\in(0,20)\). \(A'(x)=20-2x=0\Rightarrow x=10\): un cuadrado de 100 m\(^2\).
\end{ejemplo}

\begin{ojo}
\begin{itemize}
  \item Concluir que hay asíntota vertical sin comprobar que el numerador no se anula.
  \item Olvidar el dominio del logaritmo o del problema.
  \item Confundir el sentido: \(f(x-2)\) se mueve a la derecha.
  \item Dar un extremo sin comprobar el cambio de signo de \(f'\).
  \item En optimización, no comparar con los extremos del intervalo.
\end{itemize}
\end{ojo}

\end{multicols}
\end{document}
